\section{Discussion and threats to validity}\label{sec:discussion}
\subsection{Implications for operators}
Begin a sampling comparison by locating the decision and the measured cost. A gate before Collector ingress, an SDK decision, and a Collector processor avoid different work. Report accepted spans, exported bytes, Collector CPU, and sampled memory separately. Match retained IDs when comparing placement, and inspect metadata additions: equal trace selections need not yield equal serialized volume.

Treat batch formation as part of the sampling configuration. Measure exporter fanout and actual batch counts alongside timeout settings. A retain-all control is useful because it exposes buffering and grouping effects without discarding evidence. Its cost relative to bypass includes those effects, so it should not be interpreted as pure decision overhead. The results also show why a resource sign measured with JSON export should be rechecked for the deployed exporter path.

Define the diagnostic evidence before choosing a retention rate. Report protected and unprotected classes separately. For tasks using reference windows, record the number and distribution of retained reference observations as well as incident observations. Evaluate the actual scorer's full-data accuracy and an explicit deterioration margin at representative window sizes. Mean retention alone does not establish either quantity.

Finally, validate delivery independently of ideal selection. Measure decisive-span arrival relative to the first received span and decision time; test trace affinity, cache outcomes, and buffer capacity under bursts. The supporting probes in Section~\ref{sec:delivery-results} show that exported ERROR leaves can coexist with missing ancestors. Size memory for measured residence times and bursts, then verify under load.

\subsection{Construct and statistical validity}
The placement experiment controls retained trace IDs but changes ingestion, processor work, and sampling metadata together. Tail comparisons additionally change selection content and stateful processing. The fixed-input experiment controls span content while varying exporter fanout and batch timeout; actual grouping is observed. Higher-load replay also scales state capacity and adds identification attributes, so cross-campaign comparisons do not isolate load alone. Batch-size targets are fixed within campaigns, not crossed as a sensitivity factor; the RQ1 sign is established at target 1,024, not for all batch targets. These designs estimate configured treatment contrasts rather than universal per-span costs.

Resource intervals describe five or eight randomized runtime blocks on a shared node. They are pointwise, and small sample sizes do not establish the Student-$t$ assumptions. Process ticks, finite drains, and sampled RSS limit precision; sampled peaks can miss shorter transients. Replaying a fixed corpus characterizes runtime repeatability rather than independent application workloads.

The localization example is intentionally small: three HTTP handlers in one process, constructed lognormal delays, known incident boundaries, and a supplied candidate set. It measures a fixed median-change task, not unknown-incident detection or human debugging. Healthy reference windows are not no-fault incident controls. Window size changes both the full-data baseline and evidence redundancy. Calibrating on separate healthy requests prevents fault-outcome tuning of the threshold, but using the corpus-wide protected mass to equalize expected budgets remains an offline condition. A deployment must estimate that mass under changing traffic.

The 5,000 replay trials quantify uncertainty over randomized selections and incident draws from one fixed corpus under the declared iid model. They add no observed incidents. Shared priorities and nested windows create dependence between comparisons. The reported standard errors quantify conditional Monte Carlo precision and do not establish performance for particular services or strata. Inference over an incident population would require a sampling design and assumptions beyond this conditional example.

\subsection{External validity and reproducibility}
The measured resource effects apply to one shared node, loopback transport, a pinned historical Collector release, GOMAXPROCS=2, and the specified JSON or JSON-plus-Jaeger paths. JSON file export supplies an inspectable endpoint; it is not representative of every production export path. Hardware, compression, transport, trace sizes, SDKs, batching, and traffic mixtures can change resource magnitudes and signs. The study does not estimate fleet-wide savings, steady-state memory, crash durability, or long-run capacity.

Supporting delivery probes impose immediate or five-second-late arrivals relative to a two-second configured wait. They distinguish documented outcomes at those settings rather than estimating operational prevalence or a lateness-response curve. An early received root marker is a controlled input factor, not an assumption about ordinary synchronous completion order. The probability identities separately assume ideal complete-trace selection.

The study's reproducibility rests on retained observations, source freezes, configurations, and executable validation. These support inspection and reproduction but do not constitute external replication or authenticated preregistration. The primary resource comparisons and every localization configuration remain available, including excluded pilots and infeasible budgets. Alternative sampling architectures and published localizers motivate further comparative evaluation; no ranking against them follows from the measured controls.
